% ============================================================
% CHAPTER 6: LEARNING GOALS
% MCDA 5585 / 5586 Major Project Report
% Bhavik Kantilal Bhagat | A00494758 | Saint Mary's University
% ============================================================

\chapter{Learning Goals}
\label{chap:learning_goals}

The following learning goals were established at the start of my internship
in consultation with my academic and industry supervisors. They are organized
into technical and non-technical categories to reflect the dual competency
development expected of an MCDA industry placement.

% ────────────────────────────────────────────────────────────
\section{Internship Parameters}
% ────────────────────────────────────────────────────────────

My internship spanned \textbf{Mar~15 -- Aug~31, 2026}, for a total of
\textbf{900 hours} of logged work. The engagement operated in two phases:
a part-time phase from Mar~15 to Apr~15 (approximately 10--20 hours per
week) during which I completed the second semester of coursework in
parallel, followed by a full-time phase from Apr~15 to Aug~31 (approximately
40+ hours per week) upon academic semester completion. A Mitacs BSI renewal
(Sep~1 -- Dec~31, 2026) was subsequently approved, extending the industrial
collaboration into a second phase not covered by this primary submission.

I was embedded within the \textbf{Intelligent Automation and Site Reliability
Engineering (IA-SRE)} team at CTM Halifax, operating across three
interlocking work streams: SRE \& Operational Observability, Intelligent
Automation \& Predictive Operations, and Advanced Monitoring \& Capacity
Analytics.

% ────────────────────────────────────────────────────────────
\section{Technical Learning Goals}
% ────────────────────────────────────────────────────────────

\begin{itemize}

    \item \textbf{AWS Cloud Services and Infrastructure.}
          Gain production-level proficiency with the AWS services central
          to the IA-IT-SRE Infrastructure. During the part-time onboarding
          phase, I focused on understanding the full AWS service landscape
          in use at Citco: Amazon CloudWatch (metrics, alarms, Logs Insights,
          Metric Insights SQL), AWS Lambda (serverless execution, cold-start
          behaviour, concurrency), Amazon ECS and EKS (containerized
          workloads), Amazon EC2 (virtual machine infrastructure),
          Amazon MSK (managed Kafka, consumer lag), Amazon SQS
          (queue-based messaging), Amazon SNS (alerting), Amazon S3,
          AWS CloudFormation and SAM (Infrastructure-as-Code), AWS
          CodePipeline (CI/CD), and AWS Secrets Manager (credential
          management). A deliberate goal was exposure to the full spectrum
          of database categories deployed at Citco: Amazon DynamoDB
          (key-value), Amazon DocumentDB (document), Amazon Neptune
          (graph database --- CitcoWorks Knowledge Graph), Amazon RDS
          Aurora (relational), Amazon ElastiCache for Redis and Valkey
          (in-memory caching --- Meridian ODL tier), and Amazon OpenSearch
          Service (search and analytics). Additional infrastructure covered
          includes Amazon ECR, AWS CodeCommit, and the AWS IAM permission
          model (including cross-account role chaining for Grafana
          data source access).

    \item \textbf{Observability, Monitoring, and the Four Golden Signals.}
          Develop a working understanding of modern observability principles
          and apply them to a production financial services platform. This
          goal was anchored by a structured knowledge-transfer session with
          Kri Labos (SRE mentor) on the four golden signals framework:

          \begin{itemize}
              \item \textbf{Latency} --- how long requests take to complete;
                    measured at P50, P95, and P99 percentiles; degradation
                    thresholds from 200~ms to 2~s indicate performance risk.

              \item \textbf{Traffic} --- volume of demand on the system;
                    measured as request rate, Kafka throughput (messages
                    in/out), disk I/O, and Lambda invocation count; load
                    trends distinguish autoscaling candidates from capacity
                    expansion needs.

              \item \textbf{Errors} --- rate of failed requests; HTTP 5xx,
                    4xx, and 3xx error rates tracked as a percentage;
                    3--5\% threshold for SLO compliance, above 5\%
                    indicates a service or deployment issue.

              \item \textbf{Saturation} --- how full the system is; CPU,
                    memory, thread count, and database connection counts;
                    CPU at 95\% signals capacity risk; used for bottleneck
                    prediction and capacity planning.
          \end{itemize}

          Beyond the four golden signals, learning goals included
          understanding key observability metrics for each AWS service
          on the dashboard --- MSK consumer lag, SQS queue depth and
          oldest message age, ECS task CPU and memory, Lambda duration
          and concurrent executions, EC2 CPU and network, OpenSearch
          JVM heap pressure, and Redis memory utilization --- as well as
          evaluating advanced observability tooling including AWS X-Ray
          (distributed tracing), AWS Application Signals (SLO management),
          Grafana Trace Maps, and improvement opportunities for the
          existing IA ML/LLM RPA Support Agent.

    \item \textbf{Infrastructure-as-Code and CI/CD.}
          Author parameterized CloudFormation and AWS SAM templates for
          alarm and infrastructure deployment; develop a Lambda-backed
          dynamic alarm provisioner; and understand the AWS CodePipeline
          CI/CD model used for automation deployments across the platforms.

    \item \textbf{Grafana AMG Platform.}
          Build and operate a production-grade Amazon Managed Grafana
          workspace: configure CloudWatch as a data source, design
          collapsible dashboard rows for each AWS service, implement
          platform template variables for dynamic filtering, and learn
          Grafana panel types (time series, stat, bar gauge, heatmap,
          top-K list, table with deep-link URLs).

    \item \textbf{Backend API Development.}
          Write production-quality Python APIs using \texttt{Flask} and
          \texttt{FastAPI} with async route handlers, \texttt{boto3} for
          AWS service interaction, Pydantic for response validation, and
          \texttt{pytest} with Hypothesis for property-based testing.

    \item \textbf{Software Engineering Practices.}
          Practice spec-driven development via the Kiro IDE workflow
          (\texttt{.kiro/specs/} with \texttt{requirements.md},
          \texttt{design.md}, \texttt{tasks.md}), and apply AI-augmented
          engineering using Kiro (Claude Sonnet) and Amazon~Q to
          accelerate design documentation, code scaffolding, and test
          generation.

    \item \textbf{Proactive vs.\ Reactive Reliability Engineering.}
          Transition from a reactive support model (incident detection by
          business escalation) to a proactive one (alarm-backed detection
          within 30~seconds) by applying SLI/SLO targets and error budget
          concepts to real production workloads.

\end{itemize}

% ────────────────────────────────────────────────────────────
\section{Non-Technical Learning Goals}
% ────────────────────────────────────────────────────────────

\begin{itemize}

    \item \textbf{Scrum and Engineering Leadership.}
          Lead daily standup calls on rotation, facilitate team coordination,
          translate verbal requirements from Kishor into structured JIRA
          stories with clear acceptance criteria and subtask decomposition,
          and develop task estimation and prioritisation skills in a
          production engineering environment.

    \item \textbf{Cross-Functional and Cross-Timezone Collaboration.}
          Build working relationships across Halifax IA-SRE, Innovation IT
          Hyderabad (IST), Dublin and London business operations teams,
          and the RPA support tier. Participate in structured
          knowledge-transfer sessions, design reviews, security and
          compliance meetings, and stakeholder syncs across time zones.

    \item \textbf{Production Incident Response and RCA.}
          Develop the discipline of structured root cause analysis: follow
          the IA Support Model escalation tiers, document findings in SD
          tickets and Confluence, communicate clearly with business and
          platform teams under time pressure, and propose fixes or raise
          JIRA defects with actionable resolution steps.

    \item \textbf{Technical Communication and Documentation.}
          Produce engineering-grade written deliverables: incident RCA
          reports, architecture design documents, Confluence knowledge-base
          pages, runbooks, and this academic report.

    \item \textbf{Financial Domain Knowledge.}
          Develop working familiarity with fund administration operations ---
          VPL queue (5~PM~EST daily deadline), Corporate Actions workflows
          (Mandatory, Voluntary, Ops Transfer), MESO month-end cycles,
          and CAIS regulatory filing deadlines --- to make informed
          monitoring and incident response decisions.

\end{itemize}
